\chapter{Segmentación y Extracción de Línea de Costa}
\label{ch:segmentation}
% ==============================================================

\section{Segmentación semántica con SAM}

Se utiliza \textbf{Segment Anything Model (SAM)} de Meta AI en modo \textit{zero-shot}:
sin datos de entrenamiento específicos de Guardamar. A diferencia del modo "automático" de
SAM (que segmenta todo lo que encuentra en la imagen sin guía), \texttt{SAMSegmenter} usa
\textbf{\texttt{SamPredictor} con prompts de punto}: uno o varios puntos semilla por cámara
(\texttt{CAM\_PROMPTS}, coordenadas fraccionarias sobre la ROI) le dicen a SAM qué región debe
segmentar como arena seca, y se toma la máscara con mayor \texttt{score} de entre las
\texttt{multimask\_output=True} que devuelve.

\begin{codebox}[Segmentación --- extracto real de SAMSegmenter en segmentation\_sam.py]
\begin{minted}{python}
from segment_anything import sam_model_registry, SamPredictor

sam = sam_model_registry["vit_h"](checkpoint=checkpoint_path)
predictor = SamPredictor(sam)
predictor.set_image(roi_img_rgb)

# input_point/input_label vienen de CAM_PROMPTS (puntos semilla por cámara)
masks, scores, logits = predictor.predict(
    point_coords=input_point,
    point_labels=input_label,
    multimask_output=True,
)
mask = masks[np.argmax(scores)]  # la máscara con mayor score de las candidatas
\end{minted}
\end{codebox}

\begin{infobox}[¿Qué pasa si SAM no está disponible o falla?]
\texttt{analyze\_roi()} (backend/main.py) prueba una cascada de 3 niveles antes de rendirse:
\textbf{Nivel 1} SAM (si el checkpoint está instalado y no lanza excepción) →
\textbf{Nivel 2} \texttt{color\_fallback\_segmentation()} (heurística HSV, siempre disponible,
no depende de PyTorch ni de descargar el checkpoint de $\sim$2,4~GB) →
\textbf{Nivel 3} umbral de Otsu sobre el canal V, como último recurso si el color-fallback
también devuelve una máscara vacía. Cada nivel usado queda registrado en el log de actividad
(\texttt{GET /api/logs}) para poder auditar con qué método se generó cada resultado.
\end{infobox}

\section{Extracción de línea de costa}

\begin{enumerate}[leftmargin=1.5em]
  \item Cálculo de contornos sobre la máscara binaria:
        \texttt{cv2.findContours(mask, cv2.RETR\_EXTERNAL, cv2.CHAIN\_APPROX\_SIMPLE)}
  \item Selección del contorno correspondiente al borde arena húmeda -- arena seca.
  \item Resultado: polilínea en coordenadas píxel $(u, v)$.
\end{enumerate}

\section{Proyección a EPSG:25830}

Aplicación de la homografía~\eqref{eq:homography} para transformar cada punto $(u,v)$
de la polilínea a coordenadas métricas:

\begin{codebox}[pixels\_to\_utm() — proces\_images/georef\_export.py]
\begin{minted}{python}
def pixels_to_utm(points_px: list | np.ndarray, H: np.ndarray) -> np.ndarray:
    """points_px: (N,2) pares (col, row). Retorna (N,2) pares (Easting, Northing) en metros."""
    pts = np.asarray(points_px, dtype=np.float64).reshape(-1, 2)
    ones = np.ones((len(pts), 1), dtype=np.float64)
    hom = np.hstack([pts, ones])            # (N, 3)
    proj = (H @ hom.T).T                     # [X, Y, W]
    return proj[:, :2] / proj[:, 2:3]        # normalización homogénea
\end{minted}
\end{codebox}

\section{Exportación GeoJSON}

\begin{codebox}[Estructura real del GeoJSON de salida --- export\_geojson() en georef\_export.py]
\begin{minted}{json}
{
  "type": "FeatureCollection",
  "crs": {
    "type": "name",
    "properties": { "name": "urn:ogc:def:crs:EPSG::25830" }
  },
  "properties": {
    "project": "CV-LIT Guardamar",
    "created": "2026-08-11T17:00:00+00:00",
    "n_features": 1
  },
  "features": [{
    "type": "Feature",
    "geometry": {
      "type": "LineString",
      "coordinates": [[711234.5, 4209876.3], ["..."]]
    },
    "properties": {
      "ID_Camara":     "CAM_1",
      "Timestamp":     "2026-04-30T12:00:00Z",
      "Confianza_IA":  0.93,
      "Area_Seca_m2":  14250.7,
      "RMSE_m":        0.351,
      "n_puntos":      842,
      "crs":           "EPSG:25830"
    }
  }]
}
\end{minted}
\end{codebox}

\begin{infobox}[¿Por qué el CRS aparece dos veces en formatos distintos?]
El bloque \texttt{crs} de nivel superior (formato URN, el estándar GeoJSON histórico de
QGIS/OGC) identifica el sistema de referencia de la colección completa. El campo
\texttt{"crs": "EPSG:25830"} dentro de \texttt{properties} de cada feature es una copia en
formato corto pensada para lectura rápida (p.\,ej. en la tabla de \texttt{Resultados} del
frontend) sin tener que parsear el URN — ambos apuntan al mismo sistema, EPSG:25830.
\end{infobox}

% ==============================================================
